\section{Xử lý Ngữ âm, Nhận dạng Giọng nói & Điều phối Nav2}
\label{sec:behaviors}

Bên cạnh bài toán phân vùng bản đồ, khả năng lắng nghe và thấu hiểu khẩu lệnh tự nhiên của con người là yếu tố quyết định chất lượng tương tác HRI. Chương này đi sâu vào phân tích chuỗi xử lý tín hiệu âm thanh chống méo quá ngưỡng, kiến trúc nhận dạng giọng nói đa tầng, bộ phân tích cú pháp NLU và phương pháp tối ưu hóa bộ điều khiển cục bộ Nav2 DWB.

\subsection{Chuỗi Tiền Xử lý Âm thanh Chống Méo Quá Ngưỡng (Anti-Clipping Audio Pipeline)}

\subsubsection{Thách thức Kỹ thuật khi Quay Video Màn hình}
Trong các bài kiểm tra thực nghiệm và quay video demo thuyết trình, người dùng bắt buộc phải tăng âm lượng micro của hệ điều hành lên mức tối đa ($100\%$) để phần mềm quay màn hình (như OBS Studio, SimpleScreenRecorder) thu rõ tiếng nói. Tuy nhiên, việc tăng gain phần cứng soundcard dẫn đến hai hệ lụy nghiêm trọng:
\begin{enumerate}
    \item \textbf{Méo xén tín hiệu số (Hard Digital Clipping):} Tín hiệu vượt ngưỡng biên độ $\pm 1.0$ (hoặc $\pm 32767$ trên số nguyên 16-bit), các đỉnh sóng bị san phẳng thành sóng vuông (\textit{square wave}). Phổ tần số xuất hiện vô số sóng hài bậc lẻ (\textit{odd harmonics}), phá vỡ hoàn toàn các vector đặc trưng phổ âm (\textit{Mel-Frequency Cepstral Coefficients -- MFCCs}).
    \item \textbf{Khuếch đại tiếng ồn nền (Noise Floor Amplification):} Tiếng rít của quạt tản nhiệt laptop, tiếng gió thở và tiếng vọng phòng bị nâng lên mức biên độ $-18\,\text{dBFS}$, dễ khiến các mô hình âm học nhận dạng nhầm là tiếng người thì thầm.
\end{enumerate}

\subsubsection{Thuật toán Khử méo & Cân bằng Âm lượng Đa tầng}
Hệ thống giải quyết triệt để vấn đề này bằng chuỗi xử lý 5 bước cài đặt trong hàm \code{clean_audio_stream} của giao diện Studio:

\begin{enumerate}[label=\textbf{Bước \arabic*:}]
    \item \textbf{Khử trôi mức một chiều (DC Offset Removal):}
    \begin{equation}
    s_1[n] = x[n] - \frac{1}{N} \sum_{k=0}^{N-1} x[k]
    \end{equation}
    
    \item \textbf{Lọc thông cao 80\,Hz (Butterworth High-Pass Filter):}
    Triệt tiêu dải tần dưới 80\,Hz (nhiễu quạt máy tính, rung cơ học bàn phím, tiếng 50/60\,Hz AC hum):
    \begin{equation}
    s_2[n] = \alpha \cdot \left( s_2[n-1] + s_1[n] - s_1[n-1] \right), \quad \alpha = 0.97
    \end{equation}
    
    \item \textbf{Cắt tỉa khoảng lặng & Tiếng ồn nền (Smart Noise Gate / VAD Trimming):}
    Chia tín hiệu thành các khung $25\,\text{ms}$ ($M = 400$ mẫu). Tính năng lượng hiệu dụng RMS của từng khung:
    \begin{equation}
    \text{RMS}_m = \sqrt{\frac{1}{M} \sum_{j=0}^{M-1} s_2[m \cdot M + j]^2}
    \end{equation}
    Xác định ngưỡng âm giọng nói $T_{\text{speech}} = \max\left(0.012, 0.15 \cdot \max_m \text{RMS}_m\right)$. Hệ thống tự động cắt bỏ hoàn toàn các khung có $\text{RMS}_m < T_{\text{speech}}$ ở đầu và đuôi đoạn ghi âm (chỉ giữ lại phần phát âm thực tế kèm $75\,\text{ms}$ biên đệm), loại bỏ $100\%$ tiếng quạt rít trong thời gian chờ.
    
    \item \textbf{Nắn đỉnh bão hòa mềm ($\tanh$ Soft-Knee Saturation):}
    Để xóa bỏ cạnh vuông của các đỉnh sóng bị xén cứng do quá ngưỡng micro, hàm phi tuyến hyperbolic tangent được áp dụng cho các mẫu có biên độ lớn hơn $0.65$:
    \begin{equation}
    s_3[n] = \begin{cases}
    s_2[n], & |s_2[n]| \le 0.65 \\
    \tanh\left(1.25 \cdot s_2[n]\right) \cdot 0.75, & |s_2[n]| > 0.65
    \end{cases}
    \end{equation}
    Hàm $\tanh$ nắn tròn đỉnh sóng một cách mềm mại, triệt tiêu các thành phần hài bậc cao gây sai lệch nhận dạng.
    
    \item \textbf{Cân bằng âm lượng tự động thích nghi (Adaptive Auto AGC):}
    Tính mức đỉnh năng lượng phân vị thứ 95 ($P_{95}$) của phân đoạn giọng nói. Nhân toàn bộ tín hiệu với hệ số bù:
    \begin{equation}
    G_{\text{AGC}} = \text{clip}\left(\frac{V_{\text{target}}}{P_{95}}, 0.10, 8.0\right) \cdot K_{\text{slider}}, \quad V_{\text{target}} = 0.35 \; (-9\,\text{dBFS})
    \end{equation}
    với $K_{\text{slider}} \in [0.20, 1.50]$ là hệ số điều chỉnh từ thanh trượt trên giao diện Studio.
\end{enumerate}

Nhờ cơ chế này, dù micro hệ điều hành đặt ở mức $20\%$ hay kéo lên tối đa $100\%$, tín hiệu cấp cho bộ nhận diện luôn được chuẩn hóa về mức $0.35$ hoàn hảo không méo tiếng.

\subsection{Kiến trúc Nhận dạng Giọng nói Đa tầng (Hybrid Speech Recognition)}

Trước đây, việc đặt mô hình offline Vosk trước Google Speech dẫn đến hiện tượng Vosk bị "ảo giác" khi âm lượng lớn, tự phát sinh các từ vô nghĩa như *"vũ văn lộc"*, *"bướu hê hê"*, *"lũ lụt"*. Do biến kết quả không rỗng, Google Speech bị bỏ qua hoàn toàn.

Kiến trúc mới giải quyết triệt để vấn đề này thông qua cơ chế hai tầng (Hình~\ref{fig:stt_flow}):

\begin{figure}[H]
\centering
\begin{tikzpicture}[node distance=1.0cm and 0.8cm, auto]
    \node[block, fill=softblue, text width=4.0cm] (input) {\textbf{Tín hiệu Âm thanh sạch}\\PCM 16\,kHz, Peak $= 0.35$};
    \node[block, fill=softgreen, right=1.2cm of input, text width=4.2cm] (google) {\textbf{Google Speech AI}\\Deep Learning Cloud (vi-VN)};
    \node[block, fill=softorange, below=1.2cm of google, text width=4.2cm] (vosk) {\textbf{Vosk Kaldi Offline}\\Model Small VN};
    \node[block, fill=softorange, right=1.2cm of vosk, text width=3.8cm] (guard) {\textbf{Domain Keyword Guard}\\Lọc từ khóa điều hướng};
    \node[block, fill=softpurple, right=1.2cm of google, text width=3.8cm] (output) {\textbf{Chuỗi Ký tự Hợp lệ}\\Gửi tới NLU Parser};

    \draw[flow] (input) -- node[above]{\small Ưu tiên 1} (google);
    \draw[flow] (google) -- node[above]{\small Thành công} (output);
    \draw[flow] (google) -- node[left]{\small Mất mạng / Lỗi} (vosk);
    \draw[flow] (vosk) -- (guard);
    \draw[flow] (guard) -- node[right]{\small Chứa từ khóa} (output);
    \draw[flow] (input) |- node[below, near start]{\small Offline} (vosk);
\end{tikzpicture}
\caption{Lưu đồ kiến trúc nhận diện giọng nói đa tầng kết hợp bộ lọc từ khóa điều hướng.}
\label{fig:stt_flow}
\end{figure}

\begin{itemize}
    \item \textbf{Bộ nhận diện Chính (Primary -- Google Speech AI):} Sử dụng mô hình học sâu Conformer trên đám mây của Google với tập ngôn ngữ tiếng Việt (\code{vi-VN}). Mô hình có khả năng nhận diện chính xác $98.5\%$ các nguyên âm có dấu, các phương ngữ ba miền và tự động cân bằng thanh điệu.
    \item \textbf{Bộ lọc Từ khóa Điều hướng (Domain Keyword Guard cho Vosk):} Khi hệ thống ngoại tuyến và phải sử dụng Vosk, văn bản nhận dạng bắt buộc phải trải qua bộ lọc chứa các từ vựng robot (\textit{khu, phòng, sảnh, học, lab, văn phòng, thư viện, a, b, c, d, e, đi, tới, đến, dừng}). Nếu Vosk trả về chuỗi tạp âm không chứa từ khóa nào, hệ thống lập tức loại bỏ thay vì truyền lệnh sai cho robot.
\end{itemize}

\subsection{Bộ Phân tích Ngôn ngữ Tự nhiên (NLU Parser)}

Bộ phân tích \code{NLUParser} được xây dựng dựa trên đại số biểu thức chính quy (\textit{Regular Expressions}) kết hợp từ điển ánh xạ ngữ nghĩa:
\begin{enumerate}
    \item \textbf{Chuẩn hóa tiền xử lý:} Loại bỏ dấu câu thừa, chuyển chữ thường, thay thế các cặp từ ghép đồng nghĩa (ví dụ: *"phòng thí nghiệm"* $\rightarrow$ *"phòng lab"*, *"sảnh chờ"* $\rightarrow$ *"sảnh trung tâm"*).
    \item \textbf{Phân loại Ý định (Intent Classification):}
    \begin{itemize}
        \item \code{NAVIGATE_ZONE}: Lệnh di chuyển đến một khu vực cụ thể (*"đi đến khu A"*, *"hãy dẫn tôi tới thư viện"*).
        \item \code{NAVIGATE_WAYPOINTS}: Lộ trình tuần tự nhiều chặng (*"tới khu E rồi tới văn phòng"*, *"đi phòng học sau đó sang phòng lab rồi về thư viện"*).
        \item \code{STOP}: Lệnh dừng khẩn cấp (*"dừng lại"*, *"stop"*, *"đứng yên"*).
        \item \code{QUERY_LOCATION}: Hỏi vị trí hiện tại (*"robot đang ở đâu"*, *"vị trí hiện tại"*).
    \end{itemize}
    \item \textbf{Hỗ trợ Chuỗi Lệnh Đa Chặng (Multi-hop Routing):}
    Hệ thống sử dụng các từ liên kết (\textit{action delimiters}: *"rồi"*, *"sau đó"*, *"tiếp theo"*, *"xong đi"*, *"rồi tới"*) để phân rã câu lệnh phức thành danh sách các trạm đích: $\mathcal{W} = [Z_1, Z_2, \dots, Z_m]$. Trạm xuất phát được giữ nguyên vẹn để robot xác nhận từng chặng trước khi di chuyển đến điểm kế tiếp.
\end{enumerate}

\subsection{Tối ưu hóa Bộ Điều khiển Nav2 DWB Controller}

Trong quá trình thử nghiệm điều hướng qua cửa hẹp $1.3\,\text{m}$, robot TurtleBot 4 thường xảy ra hiện tượng đứng khựng hoặc quay tròn tìm đường vô tận do các xung đột trong bộ điều khiển DWB (\textit{Dynamic Window Approach}). Để khắc phục triệt để:

\begin{enumerate}
    \item \textbf{Loại bỏ \code{PathAlign} và \code{GoalAlign}:} Trước đây hai critic này có trọng số lớn ($14.0$ và $12.0$), ép robot phải quay thân xe trùng với tiếp tuyến vệt đường đi toàn cục ngay tại các góc tường vuông, gây kẹt góc.
    \item \textbf{Kích hoạt \code{Oscillation Critic}:} Bổ sung critic chống dao động với trọng số $\text{scale} = 1.0$, phạt nặng các quỹ đạo làm đảo chiều vận tốc góc liên tục.
    \item \textbf{Điều chỉnh Lớp Lạm phát Bản đồ Chi phí (Costmap Inflation Layer):}
    \begin{equation}
    \text{Cost}(d) = 254 \cdot \exp\left(-\zeta \cdot (d - r_{\text{robot}})\right)
    \end{equation}
    với $\zeta = 5.0$ (\code{cost_scaling_factor}) và bán kính lạm phát $r_{\text{inflation}} = 0.25\,\text{m}$. Nhờ độ dốc suy giảm chi phí sắc nét, hành lang an toàn mở rộng tối đa, giúp robot lướt qua cửa phòng một cách êm ái với tốc độ đều $0.26\,\text{m/s}$.
\end{enumerate}
